\documentclass[12pt,a4paper]{article}

\usepackage[french]{babel}
\usepackage[utf8]{inputenc}
\usepackage[T1]{fontenc}
\usepackage{enumerate}
\usepackage[absolute]{textpos}
\usepackage[a4paper, top=2cm, left=2cm, bottom=2cm, right=2cm]{geometry}
\usepackage{parskip}

\usepackage{graphicx}

\begin{document}
\thispagestyle{empty}

\begin{textblock*}{\paperwidth}(1.6cm,0.48cm) %{width}(horizontal, vertical)
\includegraphics[scale=0.67]{Logo_Mathematik_CMYK.pdf}
\end{textblock*}

\vspace*{2cm}

\begin{center}
\Huge{\textbf{Informations sur le déroulement de l'examen}}\\
\large{Olympiades de mathématiques 2019}
\end{center}

\bigskip

Les informations ci-dessous prévalent lors de tous les examens organisés dans le cadre des Olympiades de Mathématiques.

\section{Matériel autorisé}
\begin{itemize}
	\item Le matériel autorisé comprend des crayons/stylos ainsi que les affaires de géométrie (règle, compas, équerre).
	
	\item L'usage des calculatrices et des formulaires est interdit. Les feuilles de papier personnelles ne sont pas autorisées pendant l'examen.
	
	\item Il est possible d'apporter de la nourriture et de quoi boire pendant l'examen, mais par respect pour les autres participants nous recommandons de ne pas amener de nourriture trop odorante ou trop bruyante
\end{itemize}

\section{Déroulement de l'examen}
\begin{itemize}

	%	Les solutions seront scannées et triées par exercice pour la correction. Merci donc de ne travailler que sur un seul exercice pour chaque feuille (même pour les brouillons) et de n'utiliser que le recto des feuilles.
	
	\item Les raisonnements au brouillon et les solutions rédigées doivent être écrits uniquement sur le papier officiel fourni par les surveillants. Il est possible à tout moment de demander des feuilles supplémentaires.
	
	\item Il est possible à tout moment de poser une question relative aux énoncés des problèmes. Toutes les questions doivent être rédigées sur une feuille officielle et transmises à un surveillant
	
	%\item Tous les détails facilitant les corrections sont les bienvenus, essaye par exemple de mettre en évidence les passages importants et de marquer ce qui est superflu.
	
	%\item Il est important de correctement justifier chaque étape de ton raisonnement. Tu trouveras plus d'explications sur les bonnes pratiques en matière de rédaction de solution sur la feuille *mettre titre ici*.
	
	%\item Écris TOUT ce que tu trouves, et si tu ne sais pas comment résoudre le problème n'hésite pas à écrire toutes les observations que tu arrives à faire, même des remarques qui semblent triviales peuvent valoir des points.
	
	%\item Nous te recommandons de rendre TOUTES les feuilles lisibles, même pour des assertions complètement fausses on n'enlève jamais de points.
	
	%\item Les exercices sont triés par difficulté mais il arrive souvent que quelqu'un ne voit pas la solution d'un exercice simple tout en résolvant parfaitement un des problèmes difficiles. Ne vous laissez donc pas impressionner par le numéro d'un exercice. Tous les exercices valent le même nombre de points. Vous avez assez de temps pour vous occuper de chaque exercice et on vous conseille vivement de le faire.
	
	\item Il est possible d'aller aux toilettes pendant l'examen sur autorisation des surveillants.
	
	%\item Il est possible de rendre l'examen avant la fin et de quitter la salle sans faire de bruit.
	
\end{itemize}

\section{Rédaction et correction}
\begin{itemize}
	%\item On va t'envoyer ta copie corrigée la semaine suivant l'examen dans l'enveloppe ci-joint. S'il te plaît écris maintenant dessus ton adresse correcte et bien lisible. Si tu ne sais pas où et comment écrire l'adresse, suis l'exemple donné au tableau.

	\item Rien ne doit être écrit sur le verso des feuilles. Jamais deux exercices ne doivent être traités sur la même page.
	
	\item Sur chaque page doivent figurer le nom complet, le numéro de l'exercice, le numéro de la page et le nombre total de pages pour cet exercice. Par exemple "Augustin Cauchy, G1) p. 2/4". Chaque exercice doit être numéroté séparément.
		
	\item Toutes les étapes du raisonnement doivent être justifiées. Tous les résultats utilisés doivent être convenablement référencés.
		
	\item Lors de la réception de l'examen, des clarifications à propos de l'examen pourront être demandées dans le délai imparti. Le cas échéant, la correction pourra être réévaluée.
	
\end{itemize}

\vfill

%\begin{center}
%Bonne chance!
%\end{center}

\end{document}